\section*{Lectures on Lie Groups and Lie Algebras – 11,12}

(The material of these lectures is contained in Chapter 1 of the book
Vinberg–Onishchik, \emph{Seminar on Lie Groups and Algebraic Groups}.)

\section*{The Exponential Map}

Let $G$ be a connected Lie group and $\mathfrak g = \mathrm{Lie}\,G$.
For $x \in \mathfrak g$ denote by $\xi_x$ the corresponding right-invariant
vector field. Consider the Cauchy problem for a function
$\varphi_x : \mathbb{R} \to G$:

\[
\dot\varphi_x = \xi_x, \qquad \varphi_x(0) = e.
\]

\textbf{Proposition 1.}
The function $\varphi_x(t)$ is defined on the entire real line and satisfies

\[
\varphi_x(t+s) = \varphi_x(t)\varphi_x(s),
\qquad
\varphi_{sx}(t) = \varphi_x(st).
\]

\textbf{Proof.}
The solution exists near $0$ by the existence theorem for ODEs.
For fixed $s$, the functions $\varphi_x(t+s)$ and
$\varphi_x(t)\varphi_x(s)$ satisfy the same differential equation
$\dot\psi(t)=\xi_x$ (in the variable $t$) with the same initial condition
$\psi(0)=\varphi_x(s)$. By uniqueness, they coincide.
Hence $\varphi_x(t)$ extends to all of $\mathbb{R}$.
The second identity follows similarly.
\qed

\textbf{Definition 1.}
The exponential map $\exp : \mathfrak g \to G$ is defined by
\[
\exp(x) := \varphi_x(1).
\]

\textbf{Examples.}

1. $G=\mathbb{R}$ (additive): $\exp(x)=x$.

2. $G=\mathbb{R}_+$ (multiplicative):
\[
\exp(x)=e^x.
\]

3. $G=GL_n(\mathbb{R})$:
\[
\exp(x)=e^x.
\]

\textbf{Proposition 2.}
The map $\exp : \mathfrak g \to G$ is smooth and
\[
d_0 \exp = \mathrm{id}_{\mathfrak g}.
\]

\textbf{Proposition 3.}
The exponential map is functorial:
for any Lie group homomorphism
$\psi:G_1\to G_2$ the diagram commutes

\[
\begin{array}{ccc}
G_1 & \xrightarrow{\psi} & G_2 \\
\uparrow \exp & & \uparrow \exp \\
\mathfrak g_1 & \xrightarrow{d_e\psi} & \mathfrak g_2
\end{array}
\]

In particular:

\[
\det e^A = e^{\mathrm{tr} A},
\qquad
A^T=-A \Rightarrow e^A \in O_n.
\]

\section*{Uniqueness Theorems}

\textbf{Theorem 1.}
Any Lie algebra homomorphism
$\mathfrak g_1 \to \mathfrak g_2$
extends to a homomorphism of the corresponding connected Lie groups
in at most one way.

\textbf{Theorem 2.}
To every Lie subalgebra $\mathfrak h \subset \mathfrak g$
there corresponds at most one connected Lie subgroup
$H\subset G$ such that
\[
T_eH=\mathfrak h.
\]

\textbf{Example (Winding of the Torus).}
Let $G=S^1\times S^1$.
Choose basis $x_1,x_2$ of $\mathfrak g=\mathbb R^2$ such that

\[
\exp(ax_1+bx_2)
=
e^{2\pi i a}\times e^{2\pi i b}.
\]

Let $\mathfrak h=\mathbb R(x_1+\alpha x_2)$,
where $\alpha$ is irrational.
Then $\exp(\mathfrak h)$ is dense in $G$.
Thus $\mathfrak h$ corresponds to no Lie subgroup.

\section*{Virtual Lie Subgroups}

\textbf{Definition 2.}
A subgroup $H\subset G$ is called a \emph{virtual Lie subgroup}
if $H$ admits a Lie group structure such that
the inclusion $H\hookrightarrow G$ is a Lie group homomorphism.

For any $x\in\mathfrak g$,
\[
\{\exp(tx)\mid t\in\mathbb R\}
\]
is a virtual Lie subgroup.

\section*{Existence Theorems}

\textbf{Theorem 3.}
Every Lie algebra homomorphism
$f:\mathfrak g_1\to\mathfrak g_2$
extends to a homomorphism of connected simply connected Lie groups.

(Proof uses path lifting and ODE uniqueness.)

\textbf{Theorem 4.}
Every Lie subalgebra $\mathfrak h\subset\mathfrak g$
corresponds to a virtual Lie subgroup $H\subset G$
with
\[
T_eH=\mathfrak h.
\]

\textbf{Remark.}
By Ado's theorem,
every finite-dimensional Lie algebra admits a faithful representation
$\mathfrak h \to \mathfrak{gl}_N$.
Thus every Lie algebra is the Lie algebra of some Lie group.

\section*{Commutator and Malcev Closure}

\textbf{Proposition 4.}
The commutator subgroup $G'$ of a connected simply connected Lie group
is a normal Lie subgroup with
\[
\mathrm{Lie}\,G'=[\mathfrak g,\mathfrak g].
\]

\textbf{Definition 3.}
The Malcev closure $\mathfrak h^M$ of a Lie subalgebra
$\mathfrak h\subset\mathfrak g$
is the minimal Lie subalgebra that is tangent
to a Lie subgroup of $G$ and contains $\mathfrak h$.

\textbf{Proposition 5.}
\[
[\mathfrak h^M,\mathfrak h^M]=[\mathfrak h,\mathfrak h].
\]

\textbf{Sketch of Proof of Theorem 4.}
Reduce to the simply connected case.
Integrate $\mathfrak h^M$ and $[\mathfrak h,\mathfrak h]$
to subgroups $H^M$ and $H_0^M$.
The quotient $H^M/H_0^M$ is abelian.
Its subgroup $\exp(\mathfrak h/[\mathfrak h,\mathfrak h])$
is a virtual Lie subgroup.
Pulling it back gives the required $H$.